\documentclass[11pt,a4paper]{article}
\usepackage[T1]{fontenc}\usepackage{lmodern}
\usepackage[utf8]{inputenc}
\usepackage[margin=1in]{geometry}
\usepackage{amsmath,amssymb,microtype}
\usepackage[colorlinks=true,linkcolor=blue,urlcolor=blue]{hyperref}
\title{Loop Closure, Drift and the Interpretation of Inverse Dynamics in Golf Biomechanics}
\author{Dieter Olson}\date{Revised September 8, 2026}
\begin{document}\maketitle
\begin{abstract}
The apparent conflict between an inverse-dynamics torque plot and a golfer's sense of pushing or pulling often combines different questions. A net wrench is not a unique hand-contact distribution; a wrench moment depends on its reporting point; and a modeled input is not automatically a muscle command. None of these facts makes the moment an illusion or means that standard inverse dynamics ignores passive motion. This companion connects the wrench, contact and constrained-dynamics descriptions and identifies what additional evidence is needed for motor interpretation.
\end{abstract}

\section{Recover the Mechanical Quantity Before Interpreting It}
For a rigid club with center of mass $C$, mass $m$ and inertia $I_C$, use one inertial frame and hand-on-club action. After subtracting known non-hand loads,
\begin{align}
F_h&=ma_C-mg-F_a,\\
M_{h,C}&=I_C\alpha+\omega\times(I_C\omega)-M_{a,C}.
\end{align}
Here $a$ denotes declared non-gravitational external loading, such as aerodynamics. Uniform gravity has no moment about $C$. Ball and turf loads must be included when present. If they are unknown, the motion cannot by itself separate them from the hand wrench.

For a reference point $P$,
\begin{equation}
M_{h,P}=M_{h,C}+(r_C-r_P)\times F_h.
\end{equation}
This convention follows the usual sum of moments of forces \cite{lynch}. It does not require forces to act physically at $P$. A midpoint report is a representation of aggregate contact loading, not an assumption that all finger pressures occur at the midpoint.

The rigid approximation and measured accelerations require validation. A flexible shaft needs a model of its changing momentum and deformation. Filtering, timing, calibration and inertia errors can bias the recovered wrench. Those are scientific limitations of a particular inference; they should not be confused with the exact algebra of a reference change.

\section{Reference Geometry Does Not Create a Ghost Load}
For the same wrench at $Q$,
\begin{equation}
M_{h,Q}=M_{h,P}+(r_P-r_Q)\times F_h.
\end{equation}
A pure couple has zero resultant and is independent of origin. A general wrench's moment component changes when its nonzero resultant is represented elsewhere. An off-center force can have a real rotational effect without an independently applied local couple.

For example, a $4$ N upward force at $x=0.5$ m gives a $+2$ N m counterclockwise moment about the origin. Its equivalent representation there includes both $4$ N and $2$ N m. A pure $2$ N m couple alone is different because it has zero resultant. This is why a torque graph cannot generally tell whether the hands applied a separate push--pull couple, but also why calling its moment fictitious is wrong.

For $F=(0,F_y,0)$ and $Q-P=(s,0,0)$, the moment change is $-sF_y$ along $z$. With a chosen $F_y=40$ N and initial moment $2$ N m, moving $0.10$ m to the right gives $-2$ N m. The sign changed because the origin changed; no hand strategy has been identified. In three dimensions, the invariant $F\cdot M$ can prevent any single-force representation with zero residual moment. A presumed pressure center is not a universal solution.

\section{The Contact Map Explains a Specific Nonuniqueness}
For two hand contacts, collect their forces and local couples into $\lambda$. Their combined wrench satisfies $w_h=\mathcal G(q)\lambda$, where $\mathcal G$ is the grasp map. In force-first ordering,
\begin{equation}
F_h=F_1+F_2,\qquad
M_{h,P}=\sum_{i=1}^2[(r_i-r_P)\times F_i+C_i].
\end{equation}
At a fixed configuration, load additions in $\ker\mathcal G$ leave the rigid club wrench unchanged. This is the relevant null space; it should not be replaced by an unspecified kinematic Jacobian null space.

With point forces only, add equal-and-opposite forces \emph{along the line between the contacts}. They have zero net force and moment. Transverse opposite forces generally create a couple and are not invisible. The two-point force-only model has rank five; permitting contact couples changes the rank and attainable wrenches. Contact friction, compression and geometry determine which algebraic solutions are feasible.

The arms and torso complete the closed mechanical loop. Loop constraints restrict motion; redundant contact or actuator channels can leave load allocation underdetermined. These statements are related but not identical. An affine input model does not, by itself, imply infinitely many inputs for a specified full trajectory. Under suitable rank conditions, the net generalized input can be unique even when muscle recruitment is not.

Independent grip measurements and a validated contact model can constrain the allocation. Ground measurements add whole-body external-load information but do not automatically resolve every internal force. An optimized allocation remains conditional on its cost and constraints; it is not a recovered intention.

\section{Ordinary Inverse Dynamics Does Not Attribute All Motion to Actuation}
Use the balance
\begin{equation}
M(q)\ddot q+h(q,\dot q)=B(q)u+Q_{\rm ext}.
\end{equation}
The standard inverse output is
\begin{equation}
\tau_{\rm ID}=M\ddot q+h-Q_{\rm ext}=Bu.
\end{equation}
It already includes the declared gravity, velocity-coupling and other bias terms. Define $a_d=M^{-1}(Q_{\rm ext}-h)$; then
$\tau_{\rm ID}=M(\ddot q-a_d)$. Subtracting drift again from this output is inconsistent bookkeeping.

With $x=(q,\dot q)$, the same model is control-affine:
\begin{equation}
\dot x=f(x)+G(x)u,
\qquad f=(\dot q,a_d),\qquad G=\begin{bmatrix}0\\M^{-1}B\end{bmatrix}.
\end{equation}
The drift is defined relative to the selected input and included model terms. Calling it passive does not establish physiological relaxation. If passive tissue loading is omitted from $h$, it can appear in the estimated net load; if a constitutive model subtracts it, the residual has a different meaning.

A pendulum released under gravity has nonzero acceleration but zero pivot torque. Correct inverse dynamics returns zero for that motion. Holding the same pendulum stationary requires gravity-balancing torque. This simple check directly refutes the claim that inverse dynamics inherently assumes every departure from rest is actively caused.

\section{Constraint Reactions Can Depend on the Input}
For redundant coordinates and ideal bilateral constraints,
\begin{align}
M\ddot q+h&=Bu+Q_{\rm ext}+A^T\lambda_c,\\
A\ddot q+\dot A\dot q&=0.
\end{align}
If $A$ has independent rows, solving yields
\begin{equation}
\lambda_c=-(AM^{-1}A^T)^{-1}
[AM^{-1}(Bu+Q_{\rm ext}-h)+\dot A\dot q].
\end{equation}
The actuator term is explicit. Eliminating the reactions changes both the drift and the input gain. Zero actuator torque can coexist with nonzero reactions, but those reactions are not generally input-independent.

An ideal wrist universal joint is one possible kinematic approximation. Its constraint reactions do not prove a passive biological release; a physiological interpretation needs the permitted motion, tissue compliance, other joint action and grip conditions. An ideal workless reaction is also distinct from a dissipative or energy-storing anatomical force.

For the independent-coordinate balance, if $B$ has full column rank and the inverse load is in its image, the modeled input is uniquely determined. After eliminating constraints, use the reduced input map in this test; full rank of the original $B$ alone does not identify unknown constraint loads. With redundancy, the relevant map's kernel describes alternative allocations. If no consistent admissible allocation exists, investigate the measurements or model. Do not interpret a pseudoinverse residual as a uniquely identified muscle.

\section{Model Error, Power and Counterfactuals}
Holding measured motion fixed, omitted aerodynamics biases the inferred hand wrench by the aerodynamic wrench at the same reporting point. Its signed moment correction depends on force direction, lever arm and intrinsic aerodynamic couple. A percent change in a nearly canceled moment is not a universal swing correction.

Total rigid-body hand power is
\begin{equation}
\mathcal P_h=F_h\cdot v_P+M_{h,P}\cdot\omega.
\end{equation}
Using $v_Q=v_P+\omega\times(r_Q-r_P)$ makes the total invariant under the reference change. Its force and moment portions are separately point-dependent. A negative named torque component therefore does not necessarily mean negative total power, and mechanical power does not directly measure metabolic effort.

Similarly, an inward contact force that redirects circular motion is real even when its instantaneous power is zero. Introducing centrifugal force in a rotating frame does not create an extra physical load that can be subtracted from grip tension to recover sensation. Moving hands and a flexible shaft require the full balances rather than a fixed-pivot shortcut.

A zero-input branch follows its own subsequent state. For actual and zero-input trajectories,
\begin{equation}
\frac{d}{dt}(x-x_0)=f(x)-f(x_0)+G(x)u.
\end{equation}
Only at a common state does the drift difference disappear. Specify the intervention, horizon, output and comparison event before interpreting the difference. Setting a net torque to zero and asking a golfer to relax can produce different contact, impedance and feedback changes.

\section{A Stronger Golf and Humanoid Argument}
Earlier work shapes the state from which later dynamics evolves. Body geometry, ground contact, grip loading, shaft deformation and collision therefore interact. A measured delivery state can constrain this chain without uniquely recovering all its internal causes.

Use inverse dynamics to estimate the net load that the available evidence supports. Use contact and actuator maps to identify what remains ambiguous. Use complete work and energy balances to distinguish redirection, energy redistribution and added work. Then compare independent measurements and explicit interventions to test a physiological or coaching claim. A moment is neither a direct instruction nor an illusion; it is a precisely defined mechanical quantity whose interpretation needs the rest of the model.

\begin{thebibliography}{9}
\bibitem{lynch} K. M. Lynch and F. C. Park. \emph{Modern Robotics: Mechanics, Planning, and Control}. Cambridge University Press, 2017. Chapters 3, 8 and 12. \url{https://modernrobotics.northwestern.edu/nu-gm-book-resource/3-4-wrenches/}.
\bibitem{featherstone} R. Featherstone. \emph{Rigid Body Dynamics Algorithms}. Springer, 2008.
\end{thebibliography}
\end{document}
